\section{Dual-Loop Self-Evolving Dataset Construction}
\label{sec:dual-loop}

The workflow contains two connected loops. The inner loop operates during generation, while the outer loop operates after training and evaluation. Both loops share the same structure: an agent receives feedback from an oracle, acts within a bounded action space, and stops according to a convergence or verdict criterion. They differ mainly in time scale and target.

The inner loop is generation-time self-correction. A sample is rendered, then reviewed by VLM and CV signals for lighting, grounding, viewpoint, scale, object completeness, material plausibility, and background consistency. The controller uses this feedback to adjust Blender parameters, modify rendering scripts, improve asset placement, change camera or framing settings, filter low-quality samples, or trigger regeneration. This loop replaces part of the manual Blender tuning process that would otherwise require repeated expert inspection and code edits.

Concretely, the agent action space is deliberately practical rather than open-ended. It covers lighting and exposure adjustment, camera and framing correction, object scale and position correction, support-plane grounding, render filtering, asset regeneration, and limited script edits when repeated failures indicate that a rendering rule or controller behavior must change. These actions are bounded so that the workflow can improve samples without turning each failure into an untraceable manual rewrite.

The outer loop is deployment-time self-improvement. After a dataset round is exported, a downstream probe is trained, inference is run, metrics are computed, weak subgroups are identified, and the result is converted into an augmentation or filtering plan. That plan feeds the next data-construction round. In the current case study, this means weak-angle behavior can motivate targeted object additions, while regression guards can prevent automatic acceptance when new data introduces harmful noise.

\begin{table}[t]
\centering
\small
\caption{Inner and outer loops in the data engine.}
\label{tab:dual-loop}
\begin{tabularx}{\linewidth}{l l X X X}
\toprule
Loop & Level & Feedback oracle & Action space & Verdict role \\
\midrule
Inner loop & sample/render & VLM/CV review of generated artifacts & adjust render parameters, edit scripts, rerender, regenerate, filter & decide whether a sample or asset is usable \\
Outer loop & dataset/round & downstream LoRA probe and benchmark metrics & add targeted data, filter subsets, change quality gates, inspect or revert a round & decide whether a dataset round should continue \\
\bottomrule
\end{tabularx}
\end{table}

\begin{figure}[t]
\centering
\begin{subfigure}{0.98\linewidth}
\centering
\includegraphics[width=\linewidth]{figures/generated/inner_loop.pdf}
\caption{Generation-time self-correction at the sample/render level.}
\label{fig:inner-loop}
\end{subfigure}

\vspace{0.75em}

\begin{subfigure}{0.98\linewidth}
\centering
\includegraphics[width=\linewidth]{figures/generated/outer_loop.pdf}
\caption{Deployment-time self-improvement at the dataset/training-round level.}
\label{fig:outer-loop}
\end{subfigure}
\caption{Dual-loop self-evolving dataset construction. The inner loop replaces part of manual render tuning, while the outer loop converts downstream evaluation into the next dataset-construction decision.}
\label{fig:dual-loop}
\end{figure}

The two-loop design explains why \inspect{} is a productive outcome. A mixed result is not merely a failed experiment; it is a signal that the outer loop found a useful direction but the inner loop did not yet enforce sufficient quality. The correct response is therefore not to discard the framework, but to improve the inner quality gates before scaling further.
